\section{Well-behaved graph classes}

In this we provide a series of characteristics of the class $\cF$ which will later enable us to construct a general algorithmic framework for approximating \textsc{HC-$\cF$}. Although this characterization may occur quite technical, it is satisfied by many natural classes of graphs including trees and bounded diameter graphs (See Section \ref{subsection:FES_MMC_kBP}). The idea is that we want the problem of partitioning the graph into $\cF$-clusters in order to minimize the cost of the edges cut to admit a certain kind of natural ILP formulation with an LP relaxation solvable in polynomial time as well as a rounding algorithm for this LP with some proveable approximation guarantee. If this is the case, then it will be easy for us to encode all of the necessary constraints into our general mathemathical programming framework for approximating \textsc{HC-$\cF$}. The following definition captures the above properties in a formal way.
Any class of graphs $\cF$ that satisfies the following will be called \emph{$\alpha\br{n}$-well-behaved}:
\begin{enumerate}
    \item There exists an ILP formulation of the \textsc{P-$\cF$} problem such that: the set of variables consists only of variables of edge-cut type. An \emph{edge-cut type} variable $x_e$ for an edge $e\in E$ is an indicator of the fact that a $e$ is cut. 
    \item The property of each vertex $v\in X$ belonging to an $\cF$-cluster can be encoded using cover-type constraints (except the domain of each variable) and these are the only constraints. A \emph{cover-type} constraint is a constraint of the form $\sum_{e\in R} x_e \geq 1$ for some subset of edges $R\subseteq E$. For any $v\in X$, let $\cR_v$ denote the set of all such $R$, for which a cover-type constraint exists and is necessary for $v$ in order for it to be in an $\cF$-cluster.
    %  \DD{do dyskusji: na razie kontynuje, gdyz wiadomo o co chodzi, ale to chyba warto przeformulowac. Mowimy bowiem na poczatku, ze sa jedynie edge-cut'y, a po chwili, ze jeszcze jednak cos, a potem, ze jest wyjatek na domene.}
    \item The objective function is $\sum_{e\in E} c\br{e}\cdot x_e$.
    \item There exists a separation oracle for the LP-relaxation of the above ILP.
    \item There exists an algorithm which given an instance of the problem \textsc{P-$\cF$} and a solution to the LP relaxation of the above ILP, finds a subset of edges $S$ such that $G\setminus S$ is a valid solution to \textsc{P-$\cF$} and $\sum_{e\in S} c(e) \leq \alpha(n)\cdot \sum_{e\in E} c(e) \cdot x_e$, where $\alpha\br{n}$ is the approximation ratio of the algorithm.
\end{enumerate}

We call such a class of graphs, its ILP and its LP-relaxation \emph{$\alpha\br{n}$-well-behaved}. Although the above constraints may seem restrictive, they are satisfied by many natural classes of graphs including trees and bounded diameter graphs (See Section \ref{subsection:FES_MMC_kBP}). We encourage the reader to check that those LPs are indeed $\alpha\br{n}$-well-behaved. From now on by \textsc{P-$\cF$-ILP} and \textsc{P-$\cF$-LP} we denote any such ILP and its LP-relaxation for \textsc{P-$\cF$}. It should be noted that whenever there are multiple choices for the ILP and its LP-relaxation, one can pick arbitrarily among them, and by a slight abuse of notation the above symbols represent any of them.
